%======================================================================
\section{Light-cone wave functions and light-cone perturbation theory}
\label{sec:lcwf}
%======================================================================

\subsection{Fock space expansion with quarks}
\label{sec:fock}

The soft Fock space is now generated by the gluon, quark and antiquark creation operators acting
on the soft vacuum $\vac\equiv\vac_{\rm soft}\otimes|v\rangle_{\rm val}$ (appendix~\ref{app:states}).
A dressed state is expanded as
\begin{align}
\ket{\Psi}=\psi_{0}\vac&+\int\!d^{3}k\,\psi_{g}\ket{g}+\int\!d^{3}l\,d^{3}m\,\psi_{q\bar q}\ket{q\bar q}
\nonumber\\
&+\int\!d^{3}k\,d^{3}p\,\psi_{gg}\ket{gg}+\int\!d^{3}l\,d^{3}m\,d^{3}r\,\psi_{q\bar qg}\ket{q\bar q\,g}+\dots,
\label{eq:fock}
\end{align}
where the sums over colours, polarizations, helicities and flavours are implicit. As in Paper~I,
the expansion is in the Fock basis of the soft subspace only: we do not expand the valence part,
and the coefficients $\psi$ are functionals of the valence charge density, i.e.\ operators on
$\mathcal H_{\rm val}$. Since the $\rho$'s do not commute, Eq.~\eqref{eq:rhoalg}, their ordering
is kept throughout. Relative to the pure Yang--Mills case, the soft vacuum acquires a
quark--antiquark component at order $g^{2}$ (a gluon emitted by $\rho$ splits into a pair, or the
pair is created instantaneously off $\rho$), a soft gluon acquires a quark--antiquark component
already at order $g$, and a soft quark acquires a quark--gluon component at order $g$.

\subsection{Light-cone perturbation theory}
\label{sec:lcpt}

Light-cone perturbation theory is the Hamiltonian formulation of old fashioned perturbation
theory in light-cone quantization \cite{Lepage:1980fj,Brodsky:1997de,Kogut:1969xa}. The
interaction Hamiltonian acts on free Fock states and generates transitions between components
with different particle content; the light-cone energy of a multiparticle state is the sum of
$\eps$ over its constituents, and each intermediate state contributes an energy denominator. The
rules for fermions are the standard ones: each quark--gluon vertex contributes the matrix element
of $H_{gq\bar q}$ (appendix~\ref{app:me}), an instantaneous fermion line carrying plus momentum
$P^{+}$ contributes $1/(2P^{+})$ through $H_{ggq\bar q}$, an instantaneous gluon line contributes
$1/(P^{+})^{2}$ through $H^{(A)}_{\rm inst}$, and the fermion signs follow from the
anticommutators \eqref{eq:bb}. The derivation for non-commuting matrix elements is identical to
that of Paper~I (appendix~\ref{app:pt}).

As in Paper~I, the denominators are defined by the $i\epsilon$ prescription
$E_{i}-E_{f}\to E_{i}-E_{f}+in\epsilon$ at $n$-th order, which follows from the adiabatic definition of
$\OM$ as the M{\o}ller operator \cite{Munier:2025unitary,Kovner:2018fxj}. With
$H_{fi}\equiv\bra{f}H_{\rm int}\ket{i}$, the first and second order amplitudes are
\begin{align}
\mu_{1}(f\leftarrow i)&=\frac{H_{fi}}{E_{i}-E_{f}+i\epsilon},
\nonumber\\
\mu_{2}(f\leftarrow i)&=\sum_{m}\frac{H_{fm}H_{mi}}{(E_{i}-E_{f}+2i\epsilon)(E_{i}-E_{m}+i\epsilon)}
+\frac{H^{(2)}_{fi}}{E_{i}-E_{f}+i\epsilon},
\label{eq:mu12}
\end{align}
where $H^{(2)}$ denotes the order $g^{2}$ (instantaneous and four-point) part of the
interaction, and the ordering $H_{fm}H_{mi}$ is kept for the non-commuting coefficients. A
simple rewriting of \eqref{eq:mu12} will be used throughout this paper. Since
$H_{mi}=(E_{i}-E_{m}+i\epsilon)\,\mu_{1}(m\leftarrow i)$ and
$H_{fm}=(E_{m}-E_{f}+i\epsilon)\,\mu_{1}(f\leftarrow m)$,
\begin{equation}
\boxed{\;
\mu_{2}(f\leftarrow i)=\sum_{m}\mu_{1}(f\leftarrow m)\,\mu_{1}(m\leftarrow i)\,
\frac{E_{m}-E_{f}+i\epsilon}{E_{i}-E_{f}+2i\epsilon}
+\frac{H^{(2)}_{fi}}{E_{i}-E_{f}+i\epsilon}\;}
\label{eq:reduction}
\end{equation}
We refer to \eqref{eq:reduction} as the reduction formula: every second order amplitude is a sum,
over intermediate states, of products of first order amplitudes weighted by the fraction of the
total energy difference that is released in the second step, plus the instantaneous terms.
Two immediate consequences are worth stating.
\begin{itemize}
\item[(i)] \emph{Independent subprocesses.} If the transition $i\to f$ consists of two subprocesses
$X$ and $Y$ acting on disjoint sets of particles, with energy releases $\Delta_{X}$ and $\Delta_{Y}$,
there are two intermediate states and \eqref{eq:reduction} gives
\begin{equation}
\mu_{2}=\frac{Y X\,\Delta_{Y}+X Y\,\Delta_{X}}{\Delta_{X}+\Delta_{Y}}
\;\xrightarrow{\;[X,Y]=0\;}\;X\,Y .
\label{eq:factor}
\end{equation}
For commuting first order amplitudes the overall energy denominator cancels between the two
time orderings and the amplitude factorizes; this is the familiar partial fraction identity
$1/(D_{f}D_{X})+1/(D_{f}D_{Y})=1/(D_{X}D_{Y})$ with $D_{f}=D_{X}+D_{Y}$. Since $\rho$ appears
only in the Weizs\"acker--Williams amplitude, all the quark structures of this paper that involve
two independent subprocesses factorize exactly.
\item[(ii)] \emph{Poles.} The overall denominator $E_{i}-E_{f}$ survives in \eqref{eq:reduction} only
for transitions that can be on the energy shell, namely those that preserve the number of soft
particles of each kind ($2\to2$ scattering of soft partons, or the rescattering of a soft parton
on $\rho$). These are understood with the $i\epsilon$ prescription and their residue is the
on-shell scattering amplitude.
\end{itemize}
With the prescription \eqref{eq:mu12} the unitarity relations hold identically. In particular,
for every intermediate state $m$ \cite{PaperI},
\begin{equation}
\mu_{2}(f\leftarrow i)+\big[\mu_{2}(i\leftarrow f)\big]^{\dg}
=-\big[\mu_{1}(m\leftarrow f)\big]^{\dg}\mu_{1}(m\leftarrow i),
\label{eq:unitid}
\end{equation}
while the instantaneous terms in \eqref{eq:mu12} are anti-Hermitian and drop out of the left-hand
side. We use \eqref{eq:unitid} to verify the unitarity relations of the coefficients that are
fixed by LCPT and preserve the particle content.

\subsection{Soft wave function and the evolution operator}
\label{sec:omegadef}

As in Paper~I, the dressed eigenstates of the full soft Hamiltonian are related to the free Fock
states by the unitary evolution operator,
\begin{equation}
\ket{\Psi_{n}}=\OM\ket{n},\qquad \OM^{\dg}\OM=1,
\label{eq:Omegadef}
\end{equation}
with $\OM=U_{\epsilon}(-\infty,0)$ the adiabatic (M{\o}ller) operator \cite{Kovner:2018fxj}. At
leading order in the coupling and the highest order in $\rho$ it reduces to the coherent operator
\begin{equation}
\mathcal C=\exp\Big\{\int_{p}\Big[A^{b}_{j}(p)a^{\dg b}_{j}(p)-A^{\dg b}_{j}(p)a^{b}_{j}(p)\Big]\Big\},
\label{eq:coherent}
\end{equation}
which creates the Weizs\"acker--Williams field of the (now enlarged) valence charge. The same
operator diagonalizes the Hamiltonian: acting with $H$ on \eqref{eq:Omegadef} one finds
$\OM^{\dg}H\OM\ket{n}=E_{n}\ket{n}$, so that
\begin{equation}
H_{\rm diag}\equiv\OM^{\dg}H\OM=\sum_{n}E_{n}\ket{n}\bra{n}.
\label{eq:Hdiag}
\end{equation}
In QCD with quarks the Hamiltonian couples the gluonic Fock sectors to sectors with any number of
soft quark--antiquark pairs, and the diagonalization can only be carried out perturbatively.
